\documentclass[11pt]{article}

\usepackage[margin=1in]{geometry}
\usepackage{amsmath}
\usepackage{amssymb}
\usepackage{physics}
\usepackage{hyperref}
\usepackage{enumitem}

\title{Population Rate Equation in the Multilevel MOT Code}
\author{pMOT Monte Carlo Project}
\date{\today}

\begin{document}
\maketitle

\begin{quote}
\textbf{Archive notice.} This document describes the superseded
\texttt{mot\_error} implementation. Its use of saturated two-level scattering
expressions as elementary multilevel population rates is physically invalid;
this text is retained only for provenance.
\end{quote}

\section{What a population rate equation is}

A population rate equation is a reduced description of an atom's internal
quantum state.  Instead of tracking the full wavefunction amplitudes or the
full density matrix, it tracks only the probability that the atom occupies each
internal level.  In this project those levels are the hyperfine-Zeeman sublevels
of the \({}^{87}\mathrm{Rb}\) D2 line.  The state vector is therefore not a
position or velocity vector; it is a vector of internal-state probabilities,
\[
    \mathbf p =
    \begin{pmatrix}
    p_{g_1}, \ldots, p_{g_{N_g}},
    p_{e_1}, \ldots, p_{e_{N_e}}
    \end{pmatrix}^{T},
\]
where \(g_i\) are ground hyperfine-Zeeman states and \(e_j\) are excited
hyperfine-Zeeman states.

The central approximation is that only populations are retained.  In density
matrix language, the diagonal elements \(\rho_{ii}\) are kept and interpreted
as populations, while the off-diagonal elements \(\rho_{ij}\), \(i\ne j\), are
discarded.  Those off-diagonal elements are optical and ground-state
coherences.  They are required for phase-sensitive physics such as Rabi
oscillations, coherent dark states, coherent population trapping, Raman
coherences, standing-wave interference, and sub-Doppler polarization-gradient
cooling.  A population rate equation does not represent those effects.

The population dynamics are written as
\[
    \dv{\mathbf p}{t} = R(\mathbf r,\mathbf v)\,\mathbf p,
\]
where \(R\) is a real rate matrix.  Its off-diagonal entries are transition
rates moving population from one internal state to another.  Its diagonal
entries are negative outflow rates, chosen so that total population is
conserved.  In the current multilevel MOT implementation, the internal state is
not time-integrated during each external-motion step.  Instead, the code assumes
that optical pumping equilibrates much faster than mechanical motion and solves
the steady-state problem
\[
    R(\mathbf r,\mathbf v)\,\mathbf p_{\mathrm{ss}} = 0,
    \qquad
    \sum_i p_{\mathrm{ss},i}=1.
\]
This is the adiabatic-elimination or quasi-steady-state approximation for the
internal state.

\section{Internal basis used by the current code}

The authoritative implementation is in
\texttt{src/pmot/mot\_multilevel/rate\_equations.py} and
\texttt{src/pmot/mot\_multilevel/atomic\_structure.py}.  It uses 24 internal
states:
\begin{itemize}[nosep]
    \item 8 ground states:
    \[
        5S_{1/2}, F=1,\ m_F=-1,0,+1,
        \qquad
        5S_{1/2}, F=2,\ m_F=-2,-1,0,+1,+2.
    \]
    \item 16 excited states:
    \[
        5P_{3/2}, F'=0,1,2,3,
    \]
    with all allowed \(m_F'\) values in each manifold.
\end{itemize}

Older design text sometimes describes a 23-state model because the first
cooling-only specification omitted \(F'=0\).  The current production graph keeps
\(F'=0\), because the repumper includes the dipole-allowed
\(F=1\rightarrow F'=0\) channel.

\section{Precomputed atomic structure}

The code constructs all allowed electric-dipole transitions once in
\texttt{build\_atomic\_structure()}.  For each transition it stores:
\begin{itemize}[nosep]
    \item the ground-state index;
    \item the excited-state index;
    \item \(F,m_F\) and \(F',m_F'\);
    \item spherical polarization component \(q=m_F'-m_F\in\{-1,0,+1\}\);
    \item normalized line strength \(C^2\);
    \item excited-manifold hyperfine offset.
\end{itemize}

The transition strengths are computed from Wigner \(3j\) and \(6j\) factors and
normalized to the cycling transition
\[
    \ket{F=2,m_F=+2}
    \rightarrow
    \ket{F'=3,m_F'=+3}.
\]
The code also precomputes spontaneous-decay channels.  For every excited state,
the allowed decay branches into \(F=1\) and \(F=2\) ground states are weighted by
the same dipole strengths and normalized so that the total spontaneous decay
rate out of each excited state is \(\Gamma\).

\section{Stimulated transition rates}

At a given atom position \(\mathbf r\) and velocity \(\mathbf v\), the ordinary
multilevel MOT path evaluates the anti-Helmholtz magnetic field and uses its
direction as the local quantization axis.  Each laser beam polarization is then
decomposed into local spherical components \(\sigma^{-},\pi,\sigma^{+}\).

For each beam \(b\) and allowed transition \(g\rightarrow e\), the code computes
a transition-specific saturation parameter
\[
    s_{b e g}
    =
    \frac{I_b(\mathbf r)}{I_{\mathrm{sat}}}
    C_{eg}^{2}
    w_{bq},
\]
where \(I_b(\mathbf r)\) is the local Gaussian beam intensity,
\(C_{eg}^{2}\) is the normalized dipole strength, and \(w_{bq}\) is the local
spherical-polarization weight for the transition's required \(q\).

The effective detuning used by the code is, in angular-frequency units,
\[
    \Delta_{b e g}
    =
    \Delta_b
    - \delta_{\mathrm{HFS},eg}
    - \mathbf k_b\cdot\mathbf v
    - \mu_{eg}^{(\mathrm{Z})} |\mathbf B|/\hbar
    - \Delta^{(\mathrm{extra})}_{eg}.
\]
For the conventional multilevel MOT,
\(\Delta^{(\mathrm{extra})}_{eg}=0\).  The pMOT branch reuses the same local
environment entry point and may supply an additional transition-resonance shift.

The stimulated rate used in the current implementation is a saturated Lorentzian,
\[
    W_{b e g}
    =
    \frac{\Gamma}{2}
    \frac{s_{b e g}}
    {1+s_{b e g}+4\Delta_{b e g}^{2}/\Gamma^{2}}.
\]
The implementation treats different beams as incoherent rate contributors.  It
does not coherently sum optical fields before squaring.  This is a standard
fast MOT approximation, but it means that coherent standing-wave and
phase-sensitive saturation effects are absent.

\section{Rate-matrix construction in the current code}

The function \texttt{build\_beam\_stimulated\_rate\_matrices()} returns an array
\[
    W[b,e,g],
\]
containing the stimulated rate contributed by each beam.  Cooling beams address
ground \(F=2\) states.  Repump beams, when enabled, address ground \(F=1\)
states.  The current repumper graph permits \(F=1\rightarrow F'=0,1,2\).

The beam matrices are summed:
\[
    W_{\mathrm{tot}}[e,g] = \sum_b W[b,e,g].
\]
The full real population Liouvillian \(R\) is then assembled in
\texttt{assemble\_rate\_matrix()}:
\[
    R =
    \begin{pmatrix}
        R_{gg} & R_{ge} \\
        R_{eg} & R_{ee}
    \end{pmatrix}.
\]
The implemented block structure is:
\[
    R_{eg}[e,g] = W_{\mathrm{tot}}[e,g],
\]
\[
    R_{ge}[g,e] = W_{\mathrm{tot}}[e,g] + \Gamma_{\mathrm{decay}}[g,e].
\]
Thus the same laser-driven rate is used in reverse as a stimulated-emission
population link, and spontaneous decay is added independently.  Finally, the
diagonal entries are set column-by-column to the negative total outflow rate:
\[
    R_{ii} = -\sum_{j\ne i} R_{ji}.
\]
This enforces population conservation.

\section{Steady-state solve}

The function \texttt{steady\_state\_populations()} solves
\[
    R\mathbf p_{\mathrm{ss}}=0,
    \qquad
    \sum_i p_{\mathrm{ss},i}=1.
\]
Numerically, the code scales the matrix by \(\Gamma\), replaces one row with the
normalization condition, and calls \texttt{numpy.linalg.solve}.  If the matrix is
singular, it falls back to a least-squares solve.  Roundoff-scale negative
populations are clipped to zero, and the population vector is renormalized.

The resulting \(\mathbf p_{\mathrm{ss}}\) is interpreted as the internal-state
distribution that would be reached by fast optical pumping at the current
\((\mathbf r,\mathbf v)\).

\section{Force and diffusion}

The code computes each beam's scattering contribution by weighting that beam's
stimulated rates by the steady-state ground populations:
\[
    \Gamma_b^{(\mathrm{sc})}
    =
    \sum_{e,g}
    W[b,e,g]\,p_{\mathrm{ss}}(g).
\]
The mean optical force is then
\[
    \mathbf F
    =
    \sum_b
    \hbar \mathbf k_b\,\Gamma_b^{(\mathrm{sc})}.
\]
This is a radiation-pressure force from absorption momentum.  Gravity is added
separately during trajectory integration when enabled.  The stored force is the
optical force, not including gravity.

The code also computes a scalar recoil-diffusion coefficient,
\[
    D
    =
    \frac{1}{2}
    \sum_b
    |\hbar\mathbf k_b|^2\,
    \Gamma_b^{(\mathrm{sc})}.
\]
When trajectory diffusion is enabled, this coefficient drives a Langevin random
momentum kick.  This is an approximate recoil-heating model, not an explicit
photon-by-photon emission-direction simulation and not a full diffusion tensor.

\section{External trajectory execution}

The function \texttt{simulate\_rate\_equation\_trajectory()} advances external
motion with a fixed timestep.  At each step it:
\begin{enumerate}[nosep]
    \item evaluates the local magnetic field;
    \item chooses the local quantization axis;
    \item builds the beam-resolved stimulated rates;
    \item assembles and solves the population rate equation;
    \item computes the mean optical force and diffusion coefficient;
    \item updates mechanical momentum using optical force and gravity;
    \item optionally adds a Langevin recoil-diffusion kick;
    \item updates velocity and position;
    \item terminates if the atom escapes the configured radius.
\end{enumerate}

The trajectory update is effectively
\[
    \mathbf p_{\mathrm{mech}}(t+\Delta t)
    =
    \mathbf p_{\mathrm{mech}}(t)
    +
    \left[\mathbf F_{\mathrm{opt}}(\mathbf r,\mathbf v)
    + m\mathbf g\right]\Delta t
    +
    \sqrt{2D\Delta t}\,\boldsymbol\xi,
\]
when diffusion is enabled, where \(\boldsymbol\xi\) is a vector of independent
standard normal random variables.  Then
\[
    \mathbf v = \mathbf p_{\mathrm{mech}}/m,
    \qquad
    \mathbf r(t+\Delta t)=\mathbf r(t)+\mathbf v\Delta t.
\]

Capture and loading studies commonly disable diffusion so classifications are
deterministic and reproducible.  Temperature studies enable diffusion and
interpret the result as a finite-time Langevin temperature diagnostic.

\section{Physical fidelity and limitations}

The current population-rate model captures important MOT-scale physics:
\begin{itemize}[nosep]
    \item Rb-87 D2 hyperfine and Zeeman sublevel structure;
    \item Clebsch-Gordan-dependent coupling strengths;
    \item optical pumping among Zeeman sublevels;
    \item leakage into \(F=1\);
    \item repumper return from \(F=1\);
    \item spontaneous-decay branching;
    \item Doppler shifts;
    \item Zeeman shifts from the anti-Helmholtz field;
    \item local polarization decomposition relative to the local quantization axis;
    \item mean radiation-pressure force;
    \item approximate recoil diffusion for Langevin trajectories.
\end{itemize}

It is therefore much more realistic than a two-level MOT model for capture,
loading, force-curve, optical-pumping, and repumper-dependence studies.

However, it is not a full optical Bloch equation model.  It omits:
\begin{itemize}[nosep]
    \item optical coherences;
    \item ground-state coherences;
    \item coherent dark states;
    \item coherent population trapping;
    \item sub-Doppler and Sisyphus cooling;
    \item Rabi oscillations;
    \item coherent standing-wave interference;
    \item phase-dependent beam correlations;
    \item a full tensor recoil-diffusion model.
\end{itemize}

The model is best understood as a fast, state-resolved, Doppler-scale,
incoherent optical-pumping MOT model.  It is physically defensible when the
internal populations relax much faster than the external motion and when
coherence effects are not the dominant source of force or temperature.  It is
not sufficient for final quantitative claims about sub-Doppler temperatures,
coherent dark-state physics, or pMOT dynamics requiring a full state-resolved
AC-Stark Hamiltonian.

\section{Current status for this project}

For the conventional multilevel MOT, this rate-equation engine is the
authoritative production model, but quantitative claims remain provisional until
the relevant timestep, duration, and event-engine comparison checks are
documented for the regime being reported.

For the pMOT, the same rate-equation engine is reused as the 780-nm cooling and
repump radiation-pressure kernel.  The pMOT-specific 1529-nm trapping-light
physics is not fully represented by this population-rate equation.  The current
pMOT branch still requires a validated state-resolved AC-Stark Hamiltonian,
conservative Stark force, trapping-light scattering/heating/loss, coherent
interference treatment, and physical handling of the fictitious-field zero
before quantitative trapping claims are justified.

\end{document}
